\section{Actual interleaved family applications and collective trims}

The lemmas below construct the current-space operation at every round and
bound its viability, dimension and recorded ranks under explicit numerical
budgets. The physical filter and reference construction, the original-to-clipped
comparison, and the successful initial sample remain separate.

\begin{lemma}[One actual family image and interval trim]
    \label{lem:interval_family_trim_step}
    \lean{TNLean.MPS.AreaLaw.intervalFamilyTrimStep,
        TNLean.MPS.AreaLaw.finrank_intervalFamilyTrimStep_le,
        TNLean.MPS.AreaLaw.isIntervalSpanWithRankBound_intervalFamilyTrimStep,
        TNLean.MPS.AreaLaw.isViable_intervalFamilyTrimStep_of_tensor_decomposition}
    \leanok
    Fix an interval $[p,r)$, retaining its whole exterior. Let $T$ be a
    full-chain target and let $K=\sum_{i=1}^N F_i\otimes R_i$ in these
    interval/exterior coordinates. Suppose $K$ preserves $T$ and $T^\perp$,
    has minimum norm gain one on $T$, and squared contraction at most
    $\Delta\ge0$ on $T^\perp$. It need not be the identity on $T$.
    Let $U$ be mutually $\varepsilon$-close to $T$, with every vector of
    $U$ an MPS of bond at most $D\ge1$. Suppose $S$ is $\delta$-viable for
    $T$, where $0\le\delta<1$.

    Apply the actual family to $S$, giving $V=\sum_iF_iS$, and collectively
    trim an orthonormal basis of that actual space at threshold $\xi>0$.
    Write $S'$ for the resulting space. It has dimension at most
    $N\dim S$ and at most that many recorded spanning vectors, each with
    weighted internal-cut rank $\xi^2\operatorname{SR}(v)\le N\dim S$.
    For every $\theta>0$, put $m=\max\{r-p-1,0\}$ and
    \begin{align}
        f(\delta)&=\frac{\Delta\delta}{1-\delta+\Delta\delta},&
        E(\delta)&=(1+\theta)\bigl((1+\theta)f(\delta)
          +(1+\theta^{-1})\varepsilon+2\xi\sqrt{mD^2}\bigr)
          +(1+\theta^{-1})\varepsilon.
        \label{eq:interval_family_trim_error}
    \end{align}
    The actual space $S'$ is $E(\delta)$-viable for that same $T$.
\end{lemma}
\begin{proof}\leanok
    \uses{lem:alvv-interval-reference-weighted,lem:interval_recorded_generator_bounds}
    The exact tensor decomposition gives
    $K(S\otimes\mathcal H_{\rm out})\subseteq
      V\otimes\mathcal H_{\rm out}$.
    Normalized projection error reduction therefore makes $V$
    $f(\delta)$-viable. Apply the weighted target-relative reference
    estimate to its actual orthonormal basis and actual trim. The reference
    rank is $D^2$, since the whole exterior is retained. The image span
    has dimension at most $N\dim S$, and trimming does not increase it.
    The recorded-rank theorem applies to the actual trimmed basis vectors.
\end{proof}

\begin{lemma}[Every round of the literal interleaved iteration]
    \label{lem:interval_family_trim_iterates}
    \lean{TNLean.MPS.AreaLaw.intervalFamilyTrimIterate,
        TNLean.MPS.AreaLaw.intervalFamilyTrimIterate_zero,
        TNLean.MPS.AreaLaw.intervalFamilyTrimIterate_succ,
        TNLean.MPS.AreaLaw.finrank_intervalFamilyTrimIterate_le,
        TNLean.MPS.AreaLaw.isIntervalSpanWithRankBound_intervalFamilyTrimIterate_succ,
        TNLean.MPS.AreaLaw.isViable_intervalFamilyTrimIterate_of_tensor_decomposition}
    \leanok
    Retain one fixed target, reference, operator family and threshold as
    in Lemma~\ref{lem:interval_family_trim_step}. Starting at the actual
    initial space $S_0$, define
    \begin{align}
        S_{j+1}=\operatorname{Trim}_{\xi}
          \bigl(\operatorname{span}\{F_iS_j:i\}\bigr).
        \label{eq:interval_family_actual_iterates}
    \end{align}
    Then $\dim S_j\le N^j\dim S_0$. Every positive iterate has at most
    $N^j\dim S_0$ recorded spanning vectors, each satisfying the same
    bound on its weighted internal-cut rank. These are bounds for the
    recorded vectors, not for their entire span.

    Suppose $0\le\delta_j<1$ and $E(\delta_j)\le\delta_{j+1}$ for every
    $j$, and suppose $S_0$ is $\delta_0$-viable for $T$. Then every actual
    $S_j$ is $\delta_j$-viable for the same target.
    The recursive assignment in (\ref{eq:interval_family_actual_iterates})
    is the local correction documented in
    \cite{gap:alvv16_approximation_parameters}; it does not reapply the
    family to the unchanged initial sample.
\end{lemma}
\begin{proof}\leanok
    \uses{lem:interval_family_trim_step}
    Induct on $j$, applying the actual one-step theorem each time.
    Dimension multiplies by at most $N$; the positive-iterate recorded
    vectors are exactly the output of the last collective trim.
    Viability follows from $E(\delta_j)\le\delta_{j+1}$ and monotonicity
    in the allowed error.
\end{proof}

\begin{lemma}[Reference budgets for every actual interleaved round]
    \label{lem:interval_family_trim_reference_budget}
    \lean{TNLean.MPS.AreaLaw.isViable_intervalFamilyTrimIterate_of_reference_budgets}
    \leanok
    In Lemma~\ref{lem:interval_family_trim_iterates}, suppose $0<a\le1$,
    $\varepsilon\le a^2/16384$ and
    $\xi\sqrt{mD^2}\le a/256$. Set $\theta=a/64$. It suffices to verify
    the simpler numerical envelope
    \begin{align}
        0\le\delta_j<1,\qquad
        \frac{\Delta\delta_j}{1-\delta_j+\Delta\delta_j}+a/16
          \le\delta_{j+1}.
        \label{eq:interval_family_reference_envelope}
    \end{align}
    Every literal iterate then has the viability and dimension bounds
    of that lemma, with its actual weighted recorded ranks.

    These are target-relative reference budgets. If $T$ is a clipped
    low-energy space, its comparison with the original ground space is a
    separate accuracy condition. No fixed constant comparison error is
    silently absorbed into these small-overlap budgets.
\end{lemma}
\begin{proof}\leanok
    \uses{lem:interval_family_trim_iterates}
    Since $0\le\delta_j<1$ and $\Delta\ge0$, the normalized filtered error
    belongs to $[0,1]$. At $\theta=a/64$, the weighted reference and
    threshold terms give $E(\delta_j)\le f(\delta_j)+a/16$.
    Substitute (\ref{eq:interval_family_reference_envelope}).
\end{proof}
